\section{Discussion}

Near-ceiling classification accuracy on a strongly structured perturbation dataset
does not, by itself, distinguish between representation methods.
Logistic regression on PCA features achieves $0.981 \pm 0.003$ macro-F1 for
six-condition classification under five-fold cross-validation---essentially
saturating the linearly separable component of the task.
This ceiling arises because the MCF10A p53/CENP-A perturbations produce large,
population-scale transcriptional shifts: even with only 10\% of training labels,
logistic regression reaches 0.947 macro-F1.
In such strongly structured settings, raw classification accuracy is a poor
discriminator between representation methods.
The embedding quality analysis provides a more sensitive view: HGT-derived cell
representations achieve a silhouette score of 0.782---40-fold higher than PCA
coordinates (0.020) and 5.5-fold higher than UMAP (0.142)---indicating that the
typed graph structure organizes cells into substantially tighter condition-specific
clusters relative to inter-condition distances.
Residual classification errors, visible in the normalized confusion matrix
(Figure~\ref{fig:confusion-logistic}), concentrate at the p53-WT/DN boundary
within the chronic CENP-A overexpression condition, precisely the biologically
intermediate state where transcriptional similarity between p53 backgrounds is
highest.

scGNN \citep{wang2021scgnn} and MarsGT \citep{wu2024marsgt} report gains over
linear baselines primarily in trajectory reconstruction and rare-cell detection,
tasks where heterogeneous graph topology provides information unavailable to PCA.
DeepMAPS \citep{ma2023deepmaps} constructed cell-type-specific gene regulatory
networks from HGT, demonstrating interpretive value independent of classification
accuracy.
GEARS \citep{roohani2023gears} targets perturbation outcome prediction rather than
state classification, but shows that typed gene interaction graphs recover
combinatorial perturbation effects that linear encodings miss.
scVI and scANVI \citep{lopez2018scvi} would likely match linear classification
accuracy on this dataset given the large effect sizes, but their latent space
is optimized for reconstruction and batch correction rather than
condition-discriminative embedding.
McNemar's test confirms that the per-cell prediction patterns of logistic
regression and multitask HGT differ significantly on the held-out test split
($p<0.05$), even at comparable aggregate F1 scores, indicating that the two
models make different types of errors.
The discordant cell analysis clarifies what this difference means biologically.
Among the 13 cells that the HGT embedding classifies correctly but PCA does not,
seven are \texttt{p53-DN|acute\_oe} cells that the PCA classifier predicts as
\texttt{p53-WT|acute\_oe}.
Acute CENP-A overexpression thus attenuates the p53-status signal in the linear
PCA projection---the transcriptional states of p53-WT and p53-DN cells converge
sufficiently under acute perturbation that their first 40 principal components
overlap.
The cell k-nearest-neighbor graph incorporated in the HGT retains the
neighbourhood-level p53 signal, recovering these misassigned cells.
This finding concretely illustrates the value of graph-based representations in
perturbation settings where individual cells occupy transitional or boundary
expression states that linear projections conflate.

Several extensions follow naturally from these results.
The label-efficiency analysis shows that logistic regression on PCA reaches
near-ceiling above 30\% labeled data; graph-based neighborhood propagation
is most valuable for subtler datasets with continuous state transitions or
sparse label regimes.
The current cell-neighborhood graph uses a fixed $k$-NN topology; adaptive or
learned graph construction could better capture heterogeneous trajectory
structures \citep{wang2021scgnn}.
Gene modules are derived by k-means on expression profiles; replacing them with
curated pathway graphs from STRING or Reactome would improve biological
interpretability while preserving the typed graph structure.
Finally, the external enrichment analysis implicates translation, stress-response,
epithelial-keratin, and chromatin programs in the perturbation signature;
causal validation requires follow-up perturbational experiments beyond the
scope of this computational study.
